\documentclass[11pt]{article}
\usepackage{pictex}
\setcounter{secnumdepth}{0}
\date{1988}
\begin{document}
\title{Does Biology Constrain Culture?\thanks{Published in 
\emph{American Anthropologist}, 90(4):819--831, 1988}}
\author{Alan R. Rogers \thanks{Phone: 801/581-5529}\\
Department of Anthropology\\
University of Utah\\
Salt Lake City, UT 84112}
\maketitle
\begin{abstract}
Most social scientists would agree that the capacity for human culture was
probably fashioned by natural selection, but they disagree about the
implications of this supposition.  Some believe that natural selection
imposes important constraints on the ways in which culture can vary, while
others believe that any such constraints must be negligible.  This paper
employs a ``thought experiment'' to demonstrate that neither of these
positions can be justified by appeal to general properties of culture or
of evolution.  Natural selection can produce mechanisms of cultural
transmission that are neither adaptive nor consistent with the predictions
of \emph{acultural} evolutionary models (those ignoring cultural
evolution).  On the other hand, natural selection can also produce
mechanisms of cultural transmission that are highly consistent with
acultural models.  Thus, neither side of the sociobiology debate is
justified in dismissing the arguments of the other.  Natural selection may
impose significant constraints on some human behaviors, but negligible
constraints on others.  Models of simultaneous genetic/cultural evolution
will be useful in identifying domains in which acultural evolutionary
models are, and are not, likely to be useful. 
\end{abstract}

\section{Introduction}

The term \emph{culture} means different things to different people
(Kroeber and Kluckhohn 1952).  At least within anthropology, however,
there is consensus on one point: culture is viewed as something that we
learn from each other. When individuals learn from each other, each
generation may inherit the knowledge, attitudes, and behaviors of its
predecessors even if these are not transmitted genetically.  This
constitutes a mechanism of \emph{cultural inheritance}.  I use a definition
of culture that is broad enough to encompass most others; the term will
refer here to any pattern of behavior that is influenced by cultural
inheritance. 

This definition emphasizes that social learning is the basis of culture.  Most
social scientists would agree that human learning mechanisms were, in all
likelihood, crafted by natural selection.  They are in profound
disagreement, however, about the implications of this supposition.  Some
scientists believe that the action of natural selection on learning
mechanisms constrains culture in important ways, while others are convinced
that any such constraints are negligible.

This disagreement is most clearly displayed in the debate over human
sociobiology.  Critics of sociobiology often argue as follows:
\begin{quote}
Argument for Weak Constraints
\begin{list}{\ }{\leftmargin 2em}
\item[1.1]  Any biological constraints on culture are too weak to be
significant. 
\item[1.2] Therefore, models of genetic evolution will be of little use in
understanding variation in human behavior.
\end{list}
\end{quote}
Thus, evolutionary biology has little to offer the social sciences, and
further progress in that field will require attention to the dynamics of
cultural transmission (Science for the People 1976; Gould and Lewontin
1979; Harris 1979; Gould 1980; Kitcher 1985; Sahlins 1976).

Against this position, sociobiologists offer what Boyd and
Richerson (1985) have termed the ``argument from natural origins''
(Alexander 1979; Irons 1979; Kurland 1979; Durham 1976, 1979; Smith
1987a):
\begin{quote}
Argument from Natural Origins (i.e., For Strong Constraints)
\begin{list}{}{\leftmargin 2em}
\item[2.1] If culturally acquired behaviors tended to reduce Darwinian
fitness, then the capacity for culture would be altered or destroyed by
natural selection. 
\item[2.2] Therefore, behavior should tend to enhance fitness whether
transmission is cultural or not.
\end{list}
\end{quote}
This argument implies that biological constraints on culture are so strong
that it is reasonable for social scientists to ignore cultural
transmission.  It provides a rationale for the use of what I shall call
\emph{acultural} evolutionary models---those that ignore the dynamics of
cultural evolution---even in species where cultural transmission is known
to exist.

The argument for weak constraints has been strengthened by recent
demonstrations (Cavalli-Sforza and Feldman 1981; Boyd and Richerson 1985)
that, under some models of cultural transmission, cultural evolution leads
to patterns of behavior that are inconsistent with acultural models.
However, neither of these demonstrations is conclusive.  Cavalli-Sforza
and Feldman do not deal with the simultaneous evolution of genes and
culture, and their work therefore tells us little about the extent to
which genetic evolution constrains the evolution of culture.  Boyd and
Richerson do study simultaneous genetic/cultural evolution, but their
conclusions concerning the argument from natural origins are based on
other models in which genes and culture evolve independently.  Thus, the
extent to which culture is constrained by the natural selection is still
an open question.

A comprehensive theory of behavior, of course, would incorporate both
cultural and genetic evolution.  The dynamics of cultural evolution,
however, are complex, and the comprehensive theory is still beyond our
grasp.  For the present, all theories of behavior rely on some form
of simplification.  This article considers the merits of two such
simplifications: purely cultural models, which ignore genetic evolution,
and purely acultural models, which ignore cultural evolution.  The first
are justified only when biological constraints on culture are weak, the
second only when they are strong.

\section{Methods}

This article will evaluate the premises of the two arguments above, a much
easier task than evaluating the truth of their conclusions.  This task is
especially easy because these premises (propositions 1.1 and 2.1) have
been advanced, not as peculiar features of human culture, but as general
properties of culture or of evolution.  If they hold in general, then they
must also hold in any particular case that we care to examine.  If either
fails in \emph{any} particular case, then we can conclude that it does
not hold in general.

Rather than search for particular cases in nature, it will be convenient
to conduct a ``thought experiment''.  I shall consider the evolution of a
hypothetical species, the ``snerdwump'', which is characterized by
a rudimentary form of culture, described below.  By analysis of a
graphical and a mathematical model of simultaneous genetic and cultural
evolution, it will be shown that neither of the premises above is
justified for snerdwumps.  I will conclude that, since they fail for
snerdwumps, they cannot be general properties of culture or of evolution.

The proposition that species having culture will behave as predicted by
acultural models is often confused with another---that such species will
behave ``adaptively''.  The latter proposition is not justified by the
argument from natural origins since natural selection does not necessarily
lead to adaptation.  Apparently, this fact is appreciated mainly by
population geneticists and behavioral ecologists (Haldane 1932, p.119;
Wright 1969; Smith 1987b), for the distinction between adaptation and
agreement with acultural models is often ignored.  Thus, it will be useful
evaluate what might be called the strong form of the argument from natural
origins---the proposition that, if the capacity for culture evolved by
natural selection, then culture must be adaptive.

\section{Culture and Adaptation}

Naturalists have long been impressed with the many ways in which plants
and animals seem suited to their environments, and refer to the
characteristics that make them so as \emph{adaptations}.  The importance of
Darwin's theory stems in large part from its success in explaining these
adaptations.  This usage of ``adaptation'' is widespread in the
literature of evolutionary biology (Fisher 1958, p. 41; Haldane 1932),
and is also found in standard dictionaries.  Unfortunately,
another usage is also common.  For some evolutionists, natural
selection produces adaptation \emph{by definition} 
(Williams 1966, p. 25; Wilson and Bossert 1972, pp. 47--48; Ruyle \emph{et
al.} 1977, pp. 49).  This, however, confuses what is being explained
(adaptation) with the theory used to explain it (natural selection).  If
adaptation is defined to be the product of natural selection, then some
other word is needed to refer to the observation that individuals seem
well suited to their environments, and no such word has been suggested.
Use of the same word for both purposes is confusing, and has
probably contributed to the misconception that Darwin's theory is circular
(see Dunbar 1982).  To avoid confusion, I use the word ``adaptation'' only to
refer to the observation that organisms are somehow suited to their
environments.  In any particular case, the relationship of adaptation to
natural selection is a matter of hypothesis, not of definition. 

% -*-latex-*-
%Farmers and thieves
\begin{figure}
\begin{center}
\mbox{%
\beginpicture
\setcoordinatesystem units <3in,10in>
\setplotarea x from 0 to 1, y from 1 to 1.2
\axis left 
  label {Fitness}
  ticks withvalues {IL} {SL} / at 1.1 1.2 / /
\axis bottom 
  label {Frequency of SL}
  ticks numbered from 0 to 1 by 1 /
\axis top ticks withvalues {ESS} / at 0.556 /  /
\axis right /
% Cultural Evolution model from Rogers (1988)
%
% > wS := w + b*(1-p)*(1-u)/( 1 - p*(1-u));
%
%                                     b (1 - p) (1 - u)
%                           wS := w + -----------------
%                                       1 - p (1 - u)
% 
% > wi := w + b*(1-c);
%                                    wi := 1.1
% > wI := w + b*(1-c);
%                               wI := w + b (1 - c)
%
% Parameter values: w = b = 1, c = 0.9, s = 0, and u = 0.8.

\put {1} [t] <10pt,-6pt> at 0 1.1
\put {2} [t] <0pt,-10pt> at 0.05 1.191919192
\put {3} [b] <0pt,10pt> at 0.95 1.012345679

\plot 
0 1.2 0.05 1.191919192 0.10 1.183673469 0.15 1.175257732
0.20 1.166666667 0.25 1.157894737 0.30 1.148936170 0.35 1.139784946
0.40 1.130434783 0.45 1.120879121 0.50 1.111111111 0.55 1.101123596
0.60 1.090909091 0.65 1.080459770 0.70 1.069767442 0.75 1.058823529
0.80 1.047619048 0.85 1.036144578 0.90 1.024390244 0.95 1.012345679
1.00 1
/
% Fitness of Individual Learner
\plot 0 1.1  1 1.1 /
\setdots
\putrule from 0.556 1 to 0.556 1.2
\setsolid
\put {$\Leftarrow$} [l] at 0.556 1.15
\put {$\Rightarrow$} [r] at 0.556 1.15
\endpicture%
}
\end{center}
\caption{The fitnesses of both learning strategies and also the mean
  fitness of the population are plotted against the frequency of social
  learning, assuming that $w = b = 1$, $c = 0.9$, $s = 0$, and $u = 0.8$.  The
  numbers refer to the following assumptions: (1) the costs and benefits of
  individual learning are independent of its frequency; (2) when social
  learning is rare, its fitness value exceeds that of individual learning
  because it provides relatively up-to-date information at low cost; and (3)
  when social learning is very common, its fitness value falls below that of
  individual learning because the information it provides is out of date.
\label{fig.cult}}
\end{figure}

%%% Local Variables: 
%%% mode: latex
%%% TeX-master: "sl"
%%% End: 


Two hypotheses are involved: selection can account for adaptation only
when (1) those individuals with 
highest fitness are also those that are ``best adapted'', by the criteria
being used to recognize adaptation, and (2) selection tends to increase the
mean fitness of the population.  Fisher's (1958) ``fundamental theorem of
natural selection'' specifies circumstances under which selection will
increase mean fitness.  This theorem, however, is not nearly as general as
was once thought (Wright 1969; Ewens 1979).  When selection does not
maximize mean fitness, there is no reason to think that it will produce
adaptation.  Thus, the capacity for culture is expected to be adaptive
only if the mean fitness of a population with culture is higher than that
of one without.

Most of us would agree that human culture is an adaptation (White 1959:
8). It is widely assumed that this is to be expected if the capacity for
culture evolved through natural selection (Durham 1976).  Thus, no special
explanation has been sought for the adaptiveness of human culture.  As we
shall see, however, it is far from obvious that the forms of cultural
transmission produced by natural selection will be adaptive.

Cultural transmission occurs when individuals learn from each other as
well as from the environment.  Thus, to understand the evolution of the
capacity for culture, we must study the evolutionary forces that affect
mechanisms of learning.  The graphical argument presented below is adapted
from that of Harpending, Rogers, and Draper (1987), and incorporates ideas
developed by Boyd and Richerson (1985).  Following Boyd and Richerson, I
distinguish \emph{individual learning} (i.e., learning directly from the
environment) from \emph{social learning} (i.e., learning from others). For
example, consider the problem of deciding what foods to eat.  One
way to do this is by experiment---to eat things that you encounter, and
wait to see what happens.  If you get sick, you would be wise not to try
that food again.  Otherwise, you might include it in your diet.  This is
individual learning, and rats are very good at it.  Because of their
willingness to experiment with novel foods, they have been able to expand
over the much of the world, occupying many new habitats.  

However, individual learning entails costs as well as benefits.  Rats, for
example, undoubtedly ingest poison from time to time.  Thus, the value of
individual learning reflects a tradeoff between the benefit of a broad
diet and the cost of occasional poisoning.  A rat that learned his diet
socially, by copying an elder, would be less often poisoned, but would be
unable to use novel foods.  Thus, cultural transmission of diet reduces
the risk of being poisoned, at the cost of a narrower diet.  Similar
remarks have been made by Pulliam and Dunford (1980) and by Boyd and
Richerson (1985).

To clarify the implications of these assumptions, let us perform a
``thought experiment''.  Consider the evolution of culture in a
hypothetical species, the snerdwump, which inhabits a varying
environment.  Snerdwumps may adopt various behaviors, and the fitness
consequences of each behavior depend on the state of the environment. 
Some snerdwumps cope with this situation by individual learning---they
obtain, at some cost, information about the environment, and then exhibit
the behavior most appropriate to that environment.  Flexible responses
such as this are often favored by selection in varying environments
(Cavalli-Sforza 1974; Via and Lande 1985; Boyd and Richerson 1985).  Other
snerdwumps employ social learning; they simply adopt the behavior of some
individual---their ``cultural parent''---of the previous generation.  For
the moment, it will not matter much whether the cultural parent is also a
biological parent or not.

The implications of these assumptions are illustrated graphically in
Figure~1.  The numbers there refer to the numbered paragraphs below.
\begin{enumerate}
\item
The fitness effect of individual learning depends on its costs and
benefits, but not on what others are doing.  Thus, its fitness is a
horizontal line in Figure~1.  
\item \label{SLRare}
If social learning is very rare, nearly all cultural parents will be
individual learners.  Thus, social learners will acquire relatively recent
information about the environment.  They will acquire, that is, behaviors
that were appropriate in the immediately preceding generation.  Their
fitness will therefore exceed that of individual learners,
provided that social learning is sufficiently cheap and environmental
change sufficiently slow.
\item \label{SLCmn}
When all learning is social, no one is monitoring the environment, and
the information acquired will eventually be many generations old.
Social learning will then have lower fitness than individual learning
because, by assumption, information is worthless if sufficiently out of
date.
\end{enumerate}

These assumptions imply that the fitness of social learning exceeds that
of individual learning when its frequency is low, but falls below as its
frequency increases.  The shape of the curve connecting points
\ref{SLRare} and \ref{SLCmn}
will depend on the structure of the cultural inheritance system, i.e. on
the tendency of social learners to learn more from some individuals than
from others. The curve in Figure~1 was drawn using the mathematical
formulation presented below, but any monotonic curve connecting these
points will produce a graph of the same general form.  The conclusions
that follow hold regardless of the structure of the cultural inheritance
system.

This representation is consistent with a variety of assumptions about the
genetics of social learning.  It applies equally when social learning is
inherited as a single Mendelian character, and when what is inherited is
merely a propensity to learn socially, affected by genes at many loci.
Under most (but not all) forms of genetic inheritance, natural selection
will tend to increase the frequency of the learning mechanism with the
highest fitness, and I assume this to be the case.  From this
assumption several conclusions follow. 

The frequency of social learning among snerdwumps will increase
under natural selection until it reaches an equilibrium at the point where
the curves for individual and social learning cross, and I assume this
equilibrium to be stable.  At this equilibrium the fitness of social
learning (and the mean fitness of the population as a whole) is exactly
the same as that of individual learning; a population (or individual) with
social learning has no higher mean fitness than one without.  The
evolution of culture among snerdwumps, therefore, does not maximize
mean fitness.  
Selection often fails to maximize mean fitness when, as here, the fitness
of a morph decreases with its frequency.  This result is hardly new
(Haldane 1932; Wright 1969).  The novelty here is the
insight that the evolutionary process that produced culture is likely to
have been characterized by this form of frequency dependence.

Recall that natural selection is expected to produce adaptation only if
mean fitness is increased.  Since mean fitness is not increased in the
present model, there is no reason to expect that snerdwump culture will be
adaptive.  This does not, of course, mean that human culture is not
adaptative.  It does show, however, that if human culture is adaptive, it
must be so because of some feature not included in the present model.
Adaptation is neither an inevitable consequence of evolution nor of
culture.

Rindos (1986) agrees that evolution can produce forms of culture
that are not adaptive.  He also claims (p.316), however, that the
evolution of culture \emph{is} an adaptive process when fitness depends on
the relationship of behavior to the physical, rather than the social
environment.  The snerdwump example shows that this is not necessarily so.

These results also contradict the widely held belief that, if the capacity
for culture evolved through natural selection, then culture must, on
average, enhance fitness (Durham 1979: 43).  However, it is a mistake to
assume, as does Schwartz (1986), that sociobiology must stand or fall on
this issue, for maximization of fitness is not a general feature of
evolutionary models.  The relevance of acultural evolutionary models
depends on how well they predict human behavior, not on whether culture
maximizes fitness.  It is important to ask, then, how well such models can
be expected to predict behavior in species with cultural transmission. 

\section{How much does biology constrain culture?}

To answer a quantitative question, a quantitative model is needed,
and the one offered here is a simple formulation of the assumptions just
discussed.  Suppose that there are two environmental states, labeled 0 and
1, two behaviors labeled in the same way, and that individual fitness
depends on behavior and environment as follows:
\begin{center}
\begin{tabular}{rccc}
        &    & \multicolumn{2}{c}{Environmental State}\\
\multicolumn{2}{c}{} &   0      &      1 \\ \cline{3-4}
        & 0  & $w + b$  &   $w - b$\\
Behavior \\
        & 1  & $w - b$  &   $w + b$
\end{tabular}
\end{center}
Note that the fitness effect of each behavior is independent of the
frequency with which it is exhibited.  Thus, although the fitness effects
of snerdwump learning mechanisms are frequency-dependent, the effects
of the behaviors learned are not.  Further suppose that, in each
generation, the environment is perturbed with probability $u$, and that
after perturbations it is equally likely to be in either state.  

In the absence of any information about present or past environmental
states, the environment is equally likely to be in either state.
Therefore, the unconditional fitness of either behavior is
$w+ {1\over 2}b-{1\over 2}b=w$.
Fitness can be enhanced by obtaining information about the environment.
Suppose that the state of the environment can be determined accurately by
individual learning, but that this effort reduces fitness by $bc$ units
due to risk of predation or expenditure of time and energy.  Then the
fitness of a snerdwump adopting individual learning is
\begin{equation} \label{wI}
w_I =w+b(1-c) .
\end{equation}
How does this compare to the fitness of a snerdwump who learns socially?

The fitness, $w_S$, of a social learner, depends on the probability that
social learning will provide accurate information, and this probability
changes as cultural evolves.  Let $p$ denote the frequency after selection
of individuals adopting social learning in the current generation.  If
cultural evolution is fast compared with genetic evolution, then $p$ can
be treated as a constant for the purpose of finding what I shall call the
\emph{cultural equilibrium}, $\tilde{w}_S$.  This is the equilibrium value
of $w_S$ under the action of cultural evolution only.  Through social
learning, an individual acquires behavior that was originally acquired by
individual learning in some previous generation.  The probability that
this occurred $\tau$ generations ago is $p^{\tau - 1}(1-p)$, if
cultural parents are random members of the previous generation.  The
expected benefit obtained by adopting such behavior is zero if the
environment has been perturbed in the past $\tau$ generations, and $b$ if
no perturbation has occurred.  The probability of no perturbation is $(1 -
u)^{\tau}$.  Putting this together, we can write
\begin{eqnarray*}
\tilde{w}_{S} & = & \displaystyle w + b\sum_{\tau =1}^{\infty}
p^{\tau -1} (1 - p)(1 - u)^{\tau}  \\
     & = & \displaystyle w + b(1 - p)(1 - u)\sum_{\tau =1}^\infty
p^{\tau -1} (1 - u)^{\tau -1} .
\end{eqnarray*}
By the formula for the sum of a geometric series, this is
\[
\tilde{w}_S  = w + \frac{b(1 - p)(1 - u)}{1 - p(1 - u)} .
\]
This expression ignores the possibility that social learning may entail
costs of its own.  These could be introduced in several ways, but the
simplest results are obtained by assuming that they produce a proportional
reduction in $b$, i.e., that  
\begin{equation} \label{wS}
\tilde{w}_S  = w + \frac{b(1 - s)(1 - p)(1 - u)}
                        {1 - p(1 - u)} ,
\end{equation}
where $s$ measures the cost of social learning.  The qualitative
conclusions drawn below are unchanged if the costs in this model are made
additive instead of multiplicative.  It can be shown that $\tilde{w}_S$
is a stable equilibrium and that convergence toward it is rapid unless 
$(1 - u )p \approx 1$.  As in Figure~1, $\tilde{w}_S$ is a monotonically
decreasing function of $p$ that exceeds $w_I$ when $p$ is small, but falls
below it as $p$ increases. 

Changes in $p$, the frequency of social learning, are produced by genetic
evolution, which is assumed to be much slower than cultural evolution.  In
studying the evolution of $p$, therefore, it is reasonable to assume that
convergence toward (\ref{wS}) is effectively instantaneous.  Equation
\ref{wS} implies that $\tilde{w}_S = w_I$ when $p$ is equal to
\[
\hat{p}  = 1 - \frac{(1 - c)u}{(1 - u)(c - s)} .
\]
This is the \emph{evolutionary equilibrium} of $p$, i.e. the equilibrium
under the combined effects of cultural and genetic evolution.  It is
between zero and one provided that $(1 - u)(1 - s)  >  (1 - c)$.  

We are interested in measuring the consistency of behavior under this
model with the behavior that would be predicted by a model ignoring 
cultural transmission.  Because the fitness effects of the two behaviors
are frequency independent, an acultural model would predict that natural
selection would maximize fitness.  It would, in other words, predict
appropriate behavior in each environment.  Thus, the consistency between
behavior and the predictions of the acultural model can be measured by the
correlation between behavior and environment.  To define such a
correlation, we must first define two random variables.  Let $B$ equal
zero or unity depending on the behavior exibited by a random member of the
population, and let $S$ take the same values depending on the state of the
environment.  These variables have the same means and
variances: $E\{S\} = E\{B\} = 1/2$, and
$Var\{S\} = Var\{B\} = 1/4$.  The correlation of $S$ and $B$ is 
\begin{equation} \label{r}
r = {E\{SB\} - E\{S\}E\{B\} \over \sqrt{(Var\{S\}Var\{B\})}}
 = 4E\{SB\} - 1 ,
\end{equation}
and measures the tendency of individuals to behave appropriately.
The acultural model would predict that $r = 1$.

The value of $r$ is easy to find if $(1 - u)(1 - s) < (1 - c)$ since that
condition implies that social learning will not evolve.  Thus, all
snerdwumps are individual learners, and $r = 1$.  In the appendix it is
shown that
\begin{equation} \label{r2}
r =  {1 - c \over (1 - u)(1 - s)} ,
\end{equation}
in the more interesting case when $(1 - u)(1 - s) > (1 - c)$.  Comparison
with equation \ref{wI} reveals that the numerator here, $1 - c$, is the
proportion of the potential benefit of individual learning that is
actually realized.  It is, if you will, the efficiency of individual
learning.  Similarly, the denominator is the efficiency of social learning.
Thus, the agreement between behavior and acultural models is measured by
the ratio of the efficiencies of the two learning mechanisms.  Since both
efficiencies are positive fractions, it is clear that $r > 0$ for any set
of parameter values under which learning can evolve at all.  This means
that natural selection does constrain culture in the model considered
here.  An acultural model would explain a fraction $r^2$ of the variance
in culturally acquired behavior.  Where the efficiency of individual
learning is much lower than that of social learning, however, this
fraction may be very small, and the constraints on culture very weak.

Equation \ref{r2} can be re-expressed in terms of $p$ as follows.
\begin{equation} \label{r3}
r = \frac{1 - \hat{p}}{1 - \hat{p} (1 - u)} .
\end{equation}
This shows that $r$ is at least as large as $1 - \hat{p}$, the frequency of
individual learning.  Thus, behaviors that are acquired primarily through
individual learning should track the environment well.  Furthermore, if
environmental change is slow enough (i.e. if $u~\approx~0$), then $r$ will
be large even if most learning is social.  

However, this may require that environmental change be unrealistically
slow.  In each generation, the probability of an environmental change is
$u/2$, and this implies that the time between environmental changes is a
geometric random variable with mean $2/u$.  Using (\ref{r3}), this
quantity can be expressed as 
\begin{equation}
2/u =  {2 r \hat{p} \over (1 - r)(1 - \hat{p} )} .
\end{equation}
Note that, $r$ and $\hat p$ can both be near unity only if $2/u$ is large,
i.e., if environmental change is very slow.  For example, if
$r=\hat{p}=0.9$, then the average time between environmental changes is
$2/u=162$ generations---nearly 4,700 years for humans.  If environmental
change is sufficiently slow, selection will favor obligate behavior over
any form of learning (Cavalli-Sforza and Feldman 1983), so the model
analyzed here would be irrelevant.  Social learning cannot evolve if $2/u$
is too large, and that implies that $r$ and $\hat p$ cannot both be near
unity.  We should not expect, therefore, to find a close match between
behavior and environment when the frequency of social learning is high.

\section{Discussion}

The study of culture is in a position similar to that of biology before
Darwin.  Most of us would agree that human culture is adaptive---culturally
acquired behaviors are essential to humans, and have enabled them to
colonize habitats far outside the ecological niche of their primate
ancestors.  Yet, we have no satisfactory explanation of this adaptation.
It is possible to devise models in which natural selection produces
adaptive forms of culture (Robert Boyd, personal communication, 1987).  As we 
have seen, however, selection can also produce forms of culture that are not
adaptive.  This leaves us with an important unanswered question: What
accounts for the adaptiveness of human culture?

In answering this question, some theoretical work will be necessary.  We
must identify ways in which adaptive mechanisms of cultural inheritance
can evolve.  The model studied here ignores a variety of factors that may
tend to make culture adaptive.  For example, 
\begin{enumerate}
\item There is no cultural selection---individuals with high fitness are
no more likely than others to serve as cultural parents.
\item There are no learning biases favoring transmission of
behaviors that enhance fitness.
\item Individual learning does not modify or improve behaviors acquired
culturally.  Therefore, the value of culturally transmitted information
does not accumulate; it decays.  
\item Decisions about which form of learning to adopt are not contingent on
the quality of available environmental information.
\end{enumerate} 
Departures from any or all of these assumptions may contribute toward the
apparently adaptive nature of human culture, but their relative importance
is unclear.  Models of simultaneous genetic and cultural evolution should
help clarify this question, and useful starts in this direction have
already been made by several authors (Lumsden and Wilson 1981; Boyd and
Richerson 1985, 1988).

Many of us believe that acultural evolutionary models (those that ignore
cultural evolution) provide useful insights into the evolution of 
behavior---even that of species in which cultural transmission is
important.  One way to justify this position is what Boyd and Richerson
(1985) have called the argument from natural origins.  In its simplest
form, this argument rests on the premise that a mechanism of cultural
inheritance can evolve by natural selection only if it is adaptive.  As
we have seen, this premise is false.  However, this simple form of the
argument from natural origins has few supporters and is of no great
importance.  Most modern evolutionists are well aware that natural
selection can be maladaptive, especially where interactions between
individuals are involved.

The weak form of the argument from natural origins assumes that
mechanisms of cultural transmission can evolve only if the behaviors
produced tend to be consistent with the predictions of acultural
evolutionary models.  The consistency between behavior and acultural
models is measured here by $r$, a coefficient of correlation.  Equation
\ref{r2} shows that, at least within our hypothetical species, the
snerdwump, there is some tendency for behavior to accord with acultural
models.  However, if the efficiency of social learning is high, 
$r \approx 0$, and this tendency will be negligible.  There is therefore
no basis for the claim that, if the capacity for culture evolved by
natural selection, then culture must be explicable in terms of ordinary,
acultural evolutionary models.

On the other hand, equation \ref{r2} also shows that, if the
efficiency of individual learning rivals that of social learning,
$r \approx 1$, which implies that culture will be highly consistent with
the predictions of acultural models.  Thus, there are circumstances under
which biological constraints on culture are important.  There
is, therefore, no basis for the claim that since humans have culture,
models of genetic evolution are necessarily irrelevant to contemporary
variation in human behavior.  Thus, neither extreme position in the
sociobiology debate can be justified by appeal to general properties of
evolution or of culture.  It is possible for biological constraints on
culture to be negligible; it is also possible for them to be of
fundamental importance.  If either extreme position is correct, it must be
because of some peculiar feature of human culture that has yet to be
identified.

The advantages of theories of simultaneous genetic/cultural evolution are
obvious, but such theories are difficult to construct and it seems likely
that both purely cultural and purely acultural approaches to social
science will continue in importance.  These approaches, however, need not
arise from contradictory philosophical positions about the causes of human
behavior.  The past two decades of psychological research have undermined
the old notion that there are general laws of learning that apply
universally to all behaviors and all species (Lumsden and Wilson 1981;
Roper 1983).  Learning abilities are species specific, and are constrained
and biased in many ways.  It is likely that each mammalian species
possesses a constellation of learning abilities, each evolved for a
specific purpose.  Within a single species, biological
constraints may be strong for some behaviors, but negligible for others.
Purely acultural explanations would be needed for the former behaviors,
and purely cultural explanations for the latter.  To reconcile the two
approaches to social science, what is needed is a way to
distinguish behaviors whose evolutionary dynamics are primarily genetic
($r \approx 1$) from those whose dynamics are primarily cultural
($r \approx 0$).  This will be possible if appropriate models of
simultaneous genetic/cultural evolution can be developed.

The present model is far too simple to be advanced as a realistic
description of human culture.  The conclusions drawn so far, however, do
not depend on the model's realism but only on its consistency with what is
known about the process of natural selection.  This cannot be said of the
conclusions I am about to draw, and they should be regarded with some
skepticism.  However, experience has shown that simplistic models are
often successful in capturing the main features of phenomena that are too
complex to model in detail.  For example, the premises of the
Hardy-Weinberg theorem of population genetics almost never hold in nature,
yet its predictions are approximately correct more often than not.
Nontheless, this theorem  is
fundamental both to empirical and theoretical studies in population
genetics.  Only experience will tell whether the model advanced here is
similarly robust.  In the meantime, the model's broader implications seem
worthy of discussion.

The model assumes that the fitness effects of particular behaviors are
independent of their frequency.  This may be true of behaviors, such as
techniques of farming or foraging, with which resources are extracted from
the environment, but it is unlikely to be true of social behavior: the
success of particular social behaviors ordinarily does depend on what
others are doing.  As we have seen, such frequency dependence can
profoundly alter the outcome of evolution.  Thus, the model developed here
is probably a poor guide to the evolutionary dynamics of social behavior.

The model's main result is equation \ref{r2}, which relates the correlation
between behavior and the environment to the efficiencies of individual and
social learning.  As these undoubtedly vary from behavior to behavior, it
follows that some behaviors should track the environment well, others
poorly.  Consequently, ecology should be more useful in understanding some
behaviors than others, a fact that cultural ecologists have long
understood (Steward 1955; Oliver 1962).  This equation also implies that
acultural evolutionary models are unlikely to be helpful in understanding
behaviors that are usually learned socially. 

Behavioral ecologists have had considerable success in using acultural
models to explain aspects of human foraging.  For example, Hawkes, Hill,
and O'Connell (1982) showed that the resource species chosen by
Ach\'{e} foragers were consistent with the predictions of a simple
evolutionary model.  The availability of various food items may
change so rapidly that tradition would be a poor source of
information about the optimal choice of diet.  Thus, we might expect
a reliance on individual learning, leading to a strong correlation
between behavior and environment.  This conclusion can be obtained
from equation \ref{r2} by noting that a high rate of environmental
change (large $u$) tends to increase $r$.  It may explain the
predictive success of the acultural model used by Hawkes, Hill, and
O'Connell.  Hunting techniques, on the other hand, may be much easier
to learn socially than by trial and error ($s \ll c$), and should not
rapidly become obsolete ($u \ll 1$).  If so, equation \ref{r2}
indicates that $r \approx 0$, and acultural models should do a poor
job of explaining variation in hunting techniques. 

The results of this study are broadly consistent with the conventional
assumption that evolutionary models will best predict those behaviors with
large effects on fitness.  Such behaviors would be identified as those
with large values of $b(1 - c)$, since $b$ is the fitness benefit of
appropriate behavior and $1 - c$ the efficiency individual learning.  The
conventional assumption is partially consistent with the present model
since $r$ is also an increasing function of $1 - c$. On the other hand, $r$
depends also on the efficiency of social learning, but does not depend on
$b$.  Does this mean that acultural models may sometimes be useful even
for behaviors whose effect on fitness is small?  Perhaps, but we have
ignored the fact that each learning mechanism may affect several
behaviors.  In a more realistic model, it is likely that evolution of the
learning mechanism would be most sensitive to those behaviors with large
values of $b$.  Thus, the present results are broadly consistent with the
conventional assumption. 

\section{Conclusions}

Natural selection can produce mechanisms of cultural transmission that are
not adaptive and that may be only weakly consistent with the predictions
of \emph{acultural} evolutionary models (those that ignore cultural
evolution).  This refutes what Boyd and Richerson have called ``the
argument from natural origins'' as a general principle, and raises the
question: Why is culture adaptive?  On the other hand, natural selection
can also produce mechanisms of cultural transmission that are highly
consistent with acultural evolutionary models.  This refutes the claim
that, since humans have culture, acultural models are necessarily
irrelevant to human behavior.  It seems likely that both purely cultural
and purely acultural theories of behavior will continue in importance.
These conclusions do not depend on the realism of the model studied here. 

To the extent that the model is accurate in spite of its simplicity,
several additional conclusions can be drawn.  The correlation between
behavior and the predictions of acultural models should be better for some
behaviors than others.  It should be greatest where the efficiency of
individual learning rivals that of social learning or, equivalently, where
the frequency of social learning is low. Thus, in field studies it may be
useful to ask (1) what are the relative efficiencies of social and
individual learning in the acquisition of particular behaviors, and (2)
what are the relative frequencies of the two modes of learning?  The
dynamics of cultural evolution are likely to be most important where the
efficiency and frequency of social learning are high.  

\section*{Acknowledgements}

The name ``snerdwump'' was stolen from Henry Harpending.  
I am grateful for the comments and suggestions of R. Boyd, M. Borgerhoff
Mulder, E. Cashdan, S.  Cashdan, S. Gaulin, H. Harpending, K. Hawkes, R.
Pennington, P. Richerson, M. Salmon, R. Scaglion, M. Siegel, A. Strathern,
and E. O. Wilson.  This paper also benefited from discussion following its
presentation at the Departments of Anthropology of Emory University, the
University of Utah, and the University of Wisconsin.  The work was
supported in part by NIH grant MGN 1 R29 GM39593-01.
\appendix
\section{Appendix}
To evaluate $r$ when $(1 - u)(1 - s) > (1 - c)$, note that the term
$E\{SB\}$ in (\ref{r}) is 
\[
E\{SB\} = Prob\{\,S = 1 \mbox{ and } B = 1\,\} = 0.5 x_T ,
\]
where $x_T$ is the probability that a random individual in the
population is behaving appropriately with respect to the current
environment.  Since all individual learners adopt appropriate behavior,
$x_T$ is
\begin{equation} \label{xT}
x_T = 1 - \hat{p} (1 - x_S ) ,
\end{equation}
where $x_S$ is the probability that a social learner behaves
appropriately.  Now the fitness of a social learners is $w + (1-s)b$ with
probability $x_S$, and $w - (1-s)b$ with probability $1 - x_S$.  Thus,
\[ \hat{w}_S  = w + (1 - s) b (2 x_S  -  1 ) . \]
At equilibrium, we also have
\[ \hat{w}_S = w_I = w + b(1 - c) , \]
since both strategies have equal fitnesses at equilibrium.  Equating these
expressions leads to
\[ 1 - x_S = \frac{c - s}{2(1 - s)} , \]
and substituting this into equation \ref{xT} produces
\[ x_T = 1 - \frac{(c - s)(1 - u) - (1 - c)u}{2(1 - s)(1 - u)} . \]
Finally, substituting into the expressions for and $E\{SB\}$ and then
into equation \ref{r} produces equation \ref{r2}.
\section*{References Cited} 
\begin{description} 
\frenchspacing

\item[Alexander, R.D.]
\mbox{}\begin{description}
\item[1979] \emph{Darwinism and Human Affairs}. Seattle:
University of Washington Press.  
\end{description}

\item[Boyd, R. and P.J. Richerson]
\mbox{}\begin{description}
\item[1985] \emph{Culture and the Evolutionary Process}. Chicago:
University of Chicago Press.   
\item[1988] ``An evolutionary model of social learning: the effects of
spatial and temporal variation.'' In \emph{Social Learning: Psychological
and Biological Perspectives}, T.  Zentall and B.G. Galef, eds.
Pp.~29--48. Hillsdale, NJ:Lawrence Erlbaum, Associates. 
\end{description}

\item[Cavalli-Sforza, L.L.]
\mbox{}\begin{description}
\item[1974] The role of plasticity in biological and
cultural evolution. \emph{Annals of the New York Aca\-de\-my of Sciences}
231:43--59.
\end{description}

\item[Cavalli-Sforza, L.L and M.W. Feldman.]
\mbox{}\begin{description}
\item[1981] \emph{Cultural
Transmission and Evolution: A Quantitative Approach}. Princeton: Princeton
University Press.  
\item[1983] Cultural versus
genetic adaptation. \emph{Proceedings of the National Aca\-de\-my of
Science, USA} 80: 4993--4996.  
\end{description}

\item[Dunbar, R.I.M.]
\mbox{}\begin{description}
\item[1982] ``Adaptation, fitness, and evolutionary
tautology'', in \emph{Current Problems in Sociobiology}. King's
College Sociobiology Group, eds. Pp.~9--28. Cambridge: Cambridge
University Press. 
\end{description}

\item[Durham, W.H.]
\mbox{}\begin{description}
\item[1976] ``The adaptive significance of cultural behavior.
\emph{Human Ecology} 4: 89--121.
\item[1979] ``Natural selection, adaptation, and human social
behavior''. In \emph{Evolutionary Biology and Human Social Behavior: An
Anthropological Perspective}. N.A. Chagnon and W. Irons, eds.
Pp.~4--39. North Scituate, MA: Duxbury.
\end{description}

\item[Ewens, W.J.]
\mbox{}\begin{description}
\item[1979] \emph{Mathematical Population Genetics}.
Springer-Verlag, New York.  
\end{description}

\item[Fisher, R.A.]
\mbox{}\begin{description}
\item[1958] \emph{The Genetical Theory of Natural Selection}.
Dover, New York.  
\end{description}

\item[Gould, S.J.]
\mbox{}\begin{description}
\item[1980] ``Sociobiology and human nature: a postpanglossian
vision.''  in \emph{Sociobiology Examined}. A. Montagu, ed. 
Pp.~283--290. Oxford: Oxford University Press.   
\end{description}

\item[Gould, S.J. and R.C. Lewontin.]
\mbox{}\begin{description}
\item[1979] The spandrels of San Marco and
the Panglossian paradigm: a critique of the adaptationist programme. {\em
Proceedings of the Royal Society}, London B205:581--598.  
\end{description}

\item[Haldane, J.B.S.]
\mbox{}\begin{description}
\item[1932] \emph{The Causes of Evolution}. Cornell
University Press, Ithaca, New York.   
\end{description}

\item[Harpending, H.C., A.R. Rogers, and P. Draper.]
\mbox{}\begin{description}
\item[1987] Human
sociobiology.  \emph{Yearbook of Physical Anthropology} 30:127--150.  
\end{description}

\item[Harris, M.]
\mbox{}\begin{description}
\item[1979] \emph{Cultural Materialism: The Struggle for a
Science of Culture}.  New York: Random House. 
\end{description}

\item[Hawkes, K., K. Hill, and J.F. O'Connell.]
\mbox{}\begin{description}
\item[1982] Why hunters
gather: optimal foraging and the Ach\'{e} of eastern Paragua. {\em
American Ethnologist} 9: 379--398.  
\end{description}

\item[Irons, W.]
\mbox{}\begin{description}
\item[1979] ``Natural selection, adaptation, and human social
behavior.'' In \emph{Evolutionary Biology and Human Social Behavior: An
Anthropological Perspective}. N.A. Chagnon and W. Irons, eds.
Pp.~4--39. North Scituate, MA: Duxbury.
\end{description}

\item[Kitcher, P.]
\mbox{}\begin{description}
\item[1985] \emph{Vaulting Ambition: Sociobiology and the Quest
for Human Nature}.  Cambridge MA: MIT Press.  
\end{description}

\item[Kroeber, A.L. and C. Klukhohn.]
\mbox{}\begin{description}
\item[1952] Culture, a critical review of
the concepts and definitions. \emph{Papers of the Peabody Museum of
American Archeology and Ethnology} 47:1--223.  
\end{description}

\item[Kurland, J.]
\mbox{}\begin{description}
\item[1979] ``Paternity, mother's brother, and human
sociality.'' in \emph{Evolutionary Biology and Human Social Behavior: An
Anthropological Perspective}. N.A. Chagnon and W. Irons, eds.
Pp.~4--39. North Scituate, MA: Duxbury.
\end{description}

\item[Lumsden, C.J. and E.O. Wilson.]
\mbox{}\begin{description}
\item[1981] \emph{Genes, Mind and Culture}.
Cambridge MA: Harvard University Press. 
\end{description}

\item[Oliver, S.C.]
\mbox{}\begin{description}
\item[1962] Ecology and cultural continuity as contributing
factors in the social organization of the plains Indians.  \emph{University
of California Publications in American Archaeology and Ethnology} 48
(Reprinted 1968 in \emph{Man in Adaptation: The Cultural Present}.  Y.A.
Cohen, ed. Pp.~243--262. Chicago: Aldine).  
\end{description}

\item[Pulliam, H.R. and C. Dunford.]
\mbox{}\begin{description}
\item[1980] \emph{Programmed to Learn: An
Essay on the Evolution of Culture}. New York: Columbia University Press.  
\end{description}

\item[Rindos, David.]
\mbox{}\begin{description}
\item[1986] The evolution of the capacity for culture:
sociobiology, structuralism, and cultural selectionism. \emph{Current
Anthropology} 27:315-326.
\end{description}

\item[Roper, T.J.]
\mbox{}\begin{description}
\item[1983] Learning as a biological phenomenon, Pp.~178--212
in: (T.R. Halliday and P.J.B. Slater, Eds.) \emph{Animal Behavior. III.
Genes, Development and Learning}.  New York: W.H. Freeman.  
\end{description}

\item[Ruyle, Eugene E., F.T. Cloak, Jr., L.B. Slobodkin, and William H.
Durham.]
\mbox{}\begin{description}
\item[1977] The adaptive significance of cultural behavior: Comments and
Reply. \emph{Human Ecology} 5: 49-67.
\end{description}

\item[Sahlins, M.]
\mbox{}\begin{description}
\item[1976] \emph{The Use and Abuse of Biology: An
Anthropological Critique of Sociobiology.} Ann Arbor: University of
Michigan Press.
\end{description}

\item[Schwartz, B.]
\mbox{}\begin{description}
\item[1986] \emph{The Battle for Human Nature: Science,
Morality and Modern Life}. New York: W.W. Norton.   
\end{description}

\item[Science for the People.]
\mbox{}\begin{description}
\item[1976] Sociobiology---another biological
determinism. \emph{BioScience} 26: 182--186.
\end{description}

\item[Smith, Eric Alden.]
\mbox{}\begin{description}
\item[1987a] Folk psychology \emph{versus} pop
sociobiology. \emph{Behavioral and Brain Sciences} 10:85--86. 
\item[1987b] Optimization theory in anthropology:
applications and critiques. in \emph{The Latest on the
Best: Essays on Evolution and Optimality}, John Dupr\'e, ed.  Pp.
201--249. Cambridge, MA: MIT Press.
\end{description}

\item[Steward, J.H.]
\mbox{}\begin{description}
\item[1955] \emph{Theory of Culture Change: The Methodology
of Multilinear Evolution}. Chicago: University of Illinois Press. 
\end{description}

\item[Via, S. and R. Lande]
\mbox{}\begin{description}
\item[1985] Genotype-environment interaction and the
evolution of phenotypic plasticity. \emph{Evolution} 39:488--504.  
\end{description}

\item[White, L.A.]
\mbox{}\begin{description}
\item[1959] \emph{The Evolution of Culture: The Development of
Civilization to the Fall of Rome}. New York: McGraw-Hill.  
\end{description}

\item[Williams, G.C.]
\mbox{}\begin{description}
\item[1966] \emph{Adaptation and Natural Selection: a
Critique of Some Current Evolutionary Thought}. Princeton University
Press, Princeton, New Jersey.  
\end{description}

\item[Wilson, E.O. and W.H. Bossert.]
\mbox{}\begin{description}
\item[1972] \emph{A Primer of Population
Biology}.  Sunderland, Massachusetts: Sinauer.  
\end{description}

\item[Wright, S.]
\mbox{}\begin{description}
\item[1969] \emph{Evolution and the Genetics of Populations.
Volume 2. The Theory of Gene Frequencies}. Chicago: University of Chicago
Press.
\end{description}
\nonfrenchspacing
\end{description} 
\end{document}
